\section*{Problemas}

\begin{problema}
	\label{prob-cond-1}
	Se lanzan dos dados. Sea $A$ el evento de que la suma sea al menos 9, y sea $B$ el evento de que el primer dado muestre un número par.
	\begin{enumerate}
		\item Calcula $P(A)$
		\item Calcula $P(B)$
		\item Calcula $P(A|B)$
		\item ¿Son $A$ y $B$ independientes?
	\end{enumerate}
\end{problema}

\begin{solucion}
	El espacio muestral tiene $6 \times 6 = 36$ resultados.
	
	\begin{enumerate}
		\item $A = \{(3,6), (4,5), (4,6), (5,4), (5,5), (5,6), (6,3), (6,4), (6,5), (6,6)\}$, entonces $\card{A} = 10$ y $P(A) = \frac{10}{36} = \frac{5}{18}$.
		
		\item $B$ = primer dado es par = $\{2,4,6\} \times \{1,2,3,4,5,6\}$, entonces $\card{B} = 18$ y $P(B) = \frac{18}{36} = \frac{1}{2}$.
		
		\item $A \cap B = \{(4,5), (4,6), (6,3), (6,4), (6,5), (6,6)\}$, entonces $\card{A \cap B} = 6$ y
		\begin{align}
			P(A|B) = \frac{P(A \cap B)}{P(B)} = \frac{6/36}{18/36} = \frac{6}{18} = \frac{1}{3}
		\end{align}
		
		\item Como $P(A|B) = \frac{1}{3} \neq \frac{5}{18} = P(A)$, los eventos \textbf{no son independientes}.
	\end{enumerate}
\end{solucion}

\begin{problema}
	\label{prob-cond-2}
	Una urna contiene 5 bolas rojas y 3 bolas azules. Se extraen dos bolas sucesivamente sin reemplazo. Sea $A$ = ``la primera bola es roja'' y $B$ = ``la segunda bola es roja''.
	\begin{enumerate}
		\item Calcula $P(A)$
		\item Calcula $P(B|A)$
		\item Calcula $P(B|A')$
		\item Calcula $P(B)$ usando la regla de la probabilidad total
	\end{enumerate}
\end{problema}

\begin{solucion}
	\begin{enumerate}
		\item $P(A) = \frac{5}{8}$
		
		\item Si la primera es roja, quedan 4 rojas y 3 azules (7 bolas en total):
		\begin{align}
			P(B|A) = \frac{4}{7}
		\end{align}
		
		\item Si la primera es azul, quedan 5 rojas y 2 azules (7 bolas en total):
		\begin{align}
			P(B|A') = \frac{5}{7}
		\end{align}
		
		\item Por la regla de la probabilidad total:
		\begin{align}
			P(B) &= P(B|A)P(A) + P(B|A')P(A') \\
			&= \frac{4}{7} \cdot \frac{5}{8} + \frac{5}{7} \cdot \frac{3}{8} \\
			&= \frac{20}{56} + \frac{15}{56} = \frac{35}{56} = \frac{5}{8}
		\end{align}
	\end{enumerate}
\end{solucion}

\begin{problema}
	\label{prob-cond-3}
	En una fábrica, el 70\% de las piezas son producidas por la máquina A y el 30\% por la máquina B. La máquina A tiene una tasa de defectos del 5\%, mientras que la máquina B tiene una tasa de defectos del 10\%. Si se elige una pieza al azar:
	\begin{enumerate}
		\item ¿Cuál es la probabilidad de que la pieza sea defectuosa?
		\item Si la pieza es defectuosa, ¿cuál es la probabilidad de que provenga de la máquina A?
	\end{enumerate}
\end{problema}

\begin{solucion}
	Sea $D$ = ``la pieza es defectuosa'', $A$ = ``proviene de la máquina A'', $B$ = ``proviene de la máquina B''.
	
	\begin{enumerate}
		\item Por la regla de la probabilidad total:
		\begin{align}
			P(D) &= P(D|A)P(A) + P(D|B)P(B) \\
			&= 0.05 \times 0.70 + 0.10 \times 0.30 \\
			&= 0.035 + 0.03 = 0.065
		\end{align}
		
		\item Por el teorema de Bayes:
		\begin{align}
			P(A|D) &= \frac{P(D|A)P(A)}{P(D)} \\
			&= \frac{0.05 \times 0.70}{0.065} = \frac{0.035}{0.065} \approx 0.538
		\end{align}
	\end{enumerate}
\end{solucion}

\begin{problema}
	\label{prob-cond-4}
	Se sabe que el 80\% de los estudiantes de una universidad aprueban el examen de cálculo. De los que aprueban, el 90\% asistió a clases regularmente. De los que no aprueban, solo el 40\% asistió regularmente.
	\begin{enumerate}
		\item Si un estudiante aprobó el examen, ¿cuál es la probabilidad de que haya asistido regularmente?
		\item Si un estudiante no asistió regularmente, ¿cuál es la probabilidad de que haya aprobado?
	\end{enumerate}
\end{problema}

\begin{solucion}
	Sea $A$ = ``aprobó'', $R$ = ``asistió regularmente'', $A'$ = ``no aprobó''.
	
	\begin{enumerate}
		\item Por la regla de la probabilidad total:
		\begin{align}
			P(R|A) &= \frac{P(A|R)P(R)}{P(A)}
		\end{align}
		
		Necesitamos calcular $P(A)$:
		\begin{align}
			P(A) &= P(A|R)P(R) + P(A|R')P(R') \\
			&= 0.90 \times 0.80 + 0.40 \times 0.20 \\
			&= 0.72 + 0.08 = 0.80
		\end{align}
		
		Entonces:
		\begin{align}
			P(R|A) &= \frac{0.90 \times 0.80}{0.80} = \frac{0.72}{0.80} = 0.90
		\end{align}
		
		\item Si no asistió regularmente ($A'$), entonces $R'$:
		\begin{align}
			P(A) &= P(A|R')P(R') + P(A|R)P(R) \\
			&= 0.40 \times 0.20 + 0.90 \times 0.80 \\
			&= 0.08 + 0.72 = 0.80
		\end{align}
		
		Y por Bayes:
		\begin{align}
			P(R|A') &= \frac{P(A|R')P(R')}{P(A)} \\
			&= \frac{0.40 \times 0.20}{0.80} = \frac{0.08}{0.80} = 0.10
		\end{align}
	\end{enumerate}
\end{solucion}

\begin{problema}
	\label{prob-cond-5}
	Un fabricante produce dos tipos de piezas: tipo A con probabilidad 0.6 y tipo B con probabilidad 0.4. Las piezas tipo A tienen una tasa de defectos del 2\%, mientras que las tipo B tienen una tasa de defectos del 5\%. Si se selecciona una pieza al azar y resulta defectuosa, ¿cuál es la probabilidad de que sea tipo A?
\end{problema}

\begin{solucion}
	Sea $D$ = ``pieza defectuosa'', $A$ = ``tipo A'', $B$ = ``tipo B''.
	
	\begin{enumerate}
		\item Por la regla de la probabilidad total:
		\begin{align}
			P(D) &= P(D|A)P(A) + P(D|B)P(B) \\
			&= 0.02 \times 0.6 + 0.05 \times 0.4 \\
			&= 0.012 + 0.020 = 0.032
		\end{align}
		
		\item Por el teorema de Bayes:
		\begin{align}
			P(A|D) &= \frac{P(D|A)P(A)}{P(D)} \\
			&= \frac{0.02 \times 0.6}{0.032} = \frac{0.012}{0.032} = 0.375
		\end{align}
	\end{enumerate}
\end{solucion}

\begin{problema}
	\label{prob-cond-6}
	Demuestra que si $P(B|A) = P(B)$, entonces $A$ y $B$ son independientes.
\end{problema}

\begin{solucion}
	Si $P(B|A) = P(B)$, por definición de probabilidad condicional:
	\begin{align}
		P(B|A) &= \frac{P(A \cap B)}{P(A)} = P(B) \\
		\Rightarrow P(A \cap B) &= P(A)P(B)
	\end{align}
	
	que es precisamente la definición de independencia.
\end{solucion}
